%%% Please fill in basic information on your thesis, which will be automatically
%%% inserted at the right places. You need to replace \xxx{...} by real data.

% Type of your thesis:
%	"bc" for Bachelor's
%	"mgr" for Master's
%	"phd" for PhD
%	"rig" for rigorosum
\def\ThesisType{mgr}

% Language of your study programme:
%	"cs" for Czech
%	"en" for English
\def\StudyLanguage{cs}

% Thesis title in English (exactly as in the official assignment)
% (Note: \xxx is a "ToDo label" which makes the unfilled visible. Remove it.)
\def\ThesisTitle{Solving the diffusion equation in higher dimensions using machine learning}

% Author of the thesis (you)
\def\ThesisAuthor{Jakub Antoš}

% Year when the thesis is submitted
\def\YearSubmitted{2026}

% Name of the department or institute, where the work was officially assigned
% (according to the Organizational Structure of MFF UK in English,
% see https://www.mff.cuni.cz/en/faculty/organizational-structure,
% or a full name of a department outside MFF)
\def\Department{Institute of Theoretical Physics}

% Is it a department (katedra), or an institute (ústav)?
\def\DeptType{Institute}

% Thesis supervisor: name, surname and titles
\def\Supervisor{doc. RNDr. Michal Pavelka, Ph.D.}

% Supervisor's department (again according to Organizational structure of MFF)
\def\SupervisorsDepartment{Mathematical Institute of Charles University}

% Study programme (does not apply to rigorosum theses)
\def\StudyProgramme{Mathematical and Computational Modelling in Physics}

% An optional dedication: you can thank whomever you wish (your supervisor,
% consultant, who provided you with tea and pizza, etc.)
\def\Dedication{
I would like to thank Metacentrum for their computational resources,
my brother for his technical expertise,
MRK-consulting for their useful
advice regarding computer hardware,
and mr.framework for his patience and dedication even under
sometimes very heated conditions.
}

% Abstract (recommended length around 80-200 words; this is not a copy of your thesis assignment!)
\def\Abstract{%
This thesis explores the application of machine learning,
namely the Physics-Informed Neural Networks (PINNs),
for solving the diffusion equation in moderate to high dimensional setting.
Traditional full-grid-based approaches, such as Finite Element or Finite Difference methods,
face significant challanges in higher dimensions due to the exponential growth of computational
and memory requirements - an effect known as the curse of dimensionality.
As existing line of work adressing this is offered by
the sparse grid combination technique, which
approximates the full grid by a combination of lower-dimensional anisotropic grids
while preserving strong connection to classical solvers.
As an alternative, we explore PINNs, a reletively novel scientific machine learning paradigm
that leverages the universal approximation power of neural networks and automatic differentiation.
PINNs are mesh free methods, hence they bypass the grid construction entirely by sampling points dynamically,
offering a highly flexible alternative.
Through the numerical simulation of the general $d$-dimensional diffusion equation,
this work analyze the accuracy, computational cost, scalability, and practical limitations of both approaches.
}


% 3 to 5 keywords (recommended) separated by \sep
% Keywords are useful for indexing and searching for the theses by topic.
\def\ThesisKeywords{%
%\xxx{keyword\sep key phrase}
diffusion equations\sep
curse of dimensionality\sep
machine learning\sep
numerical methods
}

% If any of your metadata strings contains TeX macros, you need to provide
% a plain-text version for use in XMP metadata embedded in the output PDF file.
% If you are not sure, check the generated thesis.xmpdata file.
\def\ThesisAuthorXMP{\ThesisAuthor}
\def\ThesisTitleXMP{\ThesisTitle}
\def\ThesisKeywordsXMP{\ThesisKeywords}
\def\AbstractXMP{\Abstract}

% If your abstracts are long and do not fit in the infopage, you can make the
% fonts a bit smaller by this setting. (Also, you should try to compress your abstract more.)
\def\InfoPageFont{\small}
%\def\InfoPageFont{\small}  % uncomment to decrease font size

% If you are studing in a Czech programme, you also need to provide metadata in Czech:
% (in English programmes, this is not used anywhere)

\def\ThesisTitleCS{Řešení difuzní rovnice ve vyšších dimenzích pomocí strojového učení}
\def\DepartmentCS{Ústav teoretické fyziky}
\def\DeptTypeCS{Ústav}
\def\SupervisorsDepartmentCS{Matematický ústav UK}
\def\StudyProgrammeCS{Matematické a počítačové modelování ve fyzice}

\def\ThesisKeywordsCS{
difuzní rovnice\sep
prokletí dimenzionality\sep
strojové učení\sep
numerické metody
}

\def\AbstractCS{
Tato práce zkoumá využití strojového učení,
konkrétně fyzikálně informovaných neuronových sítí (PINN),
pro řešení difuzní rovnice ve středně až vysokorozměrných prostorech.
Tradiční přístupy založené na mřížkové diskretizaci,
jako je metoda konečných prvků nebo metoda konečných diferencí,
čelí ve vyšších dimenzích významným problémů kvůli exponenciálnímu nárůstu výpočetních a
paměťových nároků - jev známémý jako prokletí dimenzionality.
Jedno z existujících řešení
nabízí metoda řídkých mřížek, neboli \emph{sparse grids},
která aproximuje plnou mřížku pomocí kolekce nízkodimenzionálních anizotropních mřížek,
přičemž si zachovává silnou vazbu na klasickým numetickým metodám.
Jako alternativu k tomuto přístupu zkoumáme v této práci metody PINN, relativně nový směr ve vědeckém strojovém učení,
které využívá univerzální aproximační schopnosti neuronových sítí a automatické derivování.
Metody PINN jsou bezsíťové, což znamená, že obcházejí konstrukci mřížky díky dynamickému vzorkování bodů,
a nabízejí tak vysoce flexibilní alternativu.
Pomocí numerických simulací obecné $d$-rozměrné rovnice difúze tato práce analyzuje přesnost,
výpočetní náročnost, škálovatelnost, a praktická omezení obou těchto přístupů.
}
